%-------------------------
% Resume in Latex
% Author : Stéphane KPOVIESSI
% Target: Big Data & AI Engineer / Data Scientist / GenAI Engineer
% Format: Two-Page Resume
% Version: FR (version française maîtresse)
%-------------------------

\documentclass[letterpaper,11pt]{article}

\usepackage{latexsym}
\usepackage[empty]{fullpage}
\usepackage{titlesec}
\usepackage[usenames,dvipsnames]{color}
\usepackage{enumitem}
\usepackage[colorlinks=true, urlcolor=primary, linkcolor=primary]{hyperref}
\usepackage{fancyhdr}
\usepackage[french]{babel} 
\usepackage{tabularx}
\usepackage{needspace}
\usepackage{graphicx}
\usepackage{tikz}
\input{glyphtounicode}

% --- COULEURS ---
\definecolor{primary}{RGB}{0, 79, 144}
\definecolor{darkgrey}{RGB}{50, 50, 50}

%----------FONT OPTIONS----------
\usepackage[sfdefault]{noto-sans}
\usepackage[T1]{fontenc}

\pagestyle{fancy}
\fancyhf{}
\fancyfoot{}
\renewcommand{\headrulewidth}{0pt}
\renewcommand{\footrulewidth}{0pt}
\setlength{\footskip}{4.1pt}

% --- MARGES ---
\addtolength{\oddsidemargin}{-0.5in}
\addtolength{\evensidemargin}{-0.5in}
\addtolength{\textwidth}{1.0in}
\addtolength{\topmargin}{-0.65in}
\addtolength{\textheight}{1.45in}

\urlstyle{same}
\raggedbottom
\raggedright
\setlength{\tabcolsep}{0in}

% --- FORMATAGE ---
\titleformat{\section}{
  \vspace{-4pt}\scshape\raggedright\large\color{primary}
}{}{0em}{}[\color{primary}\titlerule \vspace{0pt}]

\pdfgentounicode=1

%-------------------------
% Custom commands
\newcommand{\resumeItem}[1]{
  \item\small{{#1 \vspace{-1.5pt}}}
}

\newcommand{\resumeSubheading}[4]{
  \vspace{1pt}\item
    \begin{tabular*}{0.97\textwidth}[t]{l@{\extracolsep{\fill}}r}
      {\color{primary}\textbf{#1}} & #2 \\
      \textit{\small#3} & \textit{\small #4} \\
    \end{tabular*}\vspace{-4pt}
}

\newcommand{\resumeProjectHeading}[2]{
    \Needspace{6\baselineskip}
    \vspace{2pt}\item
    \begin{tabular*}{0.97\textwidth}{l@{\extracolsep{\fill}}r}
      \small#1 & #2 \\
    \end{tabular*}\vspace{-3pt}
}

\newcommand{\resumeSubItem}[1]{\resumeItem{#1}\vspace{-4pt}}
\renewcommand\labelitemii{$\vcenter{\hbox{\tiny$\bullet$}}$}
\newcommand{\resumeSubHeadingListStart}{\begin{itemize}[leftmargin=0.15in, label={}]}
\newcommand{\resumeSubHeadingListEnd}{\end{itemize}}
\newcommand{\resumeItemListStart}{\begin{itemize}[leftmargin=0.15in, label={\small\textbullet}, itemsep=1.5pt, parsep=0pt, topsep=2pt, partopsep=0pt]}
\newcommand{\resumeItemListEnd}{\end{itemize}\vspace{-1pt}}

%-------------------------------------------
\begin{document}

%----------HEADING----------
\noindent
\begin{minipage}[c]{0.20\textwidth}
  \centering
  \begin{tikzpicture}
    \begin{scope}
      \clip (0,0) circle (1.45cm);
      \node at (0,0) {\includegraphics[width=2.9cm]{./profile.jpeg}};
    \end{scope}
    \draw[color=primary, line width=1.2pt] (0,0) circle (1.45cm);
  \end{tikzpicture}
\end{minipage}%
\hspace{0.01\textwidth}%
\begin{minipage}[c]{0.77\textwidth}
  \centering
  {\Huge \scshape \color{primary} Stéphane KPOVIESSI} \\[3pt]
  {\large \textbf{AI Engineer - Agents \& AI Systems}} \\[7pt]
  \small
  Île-de-France (Cergy)
    \enspace$|$\enspace Disponible en France et à l'international
    \enspace$|$\enspace +33\,7\,48\,59\,47\,13 \\[2pt]
  \href{mailto:oastephanekpoviessi@gmail.com}{oastephanekpoviessi@gmail.com}
    \enspace$|$\enspace
  \href{https://linkedin.com/in/stephanekpoviessi}{LinkedIn}
    \enspace$|$\enspace
  \href{https://github.com/EngineerProjects}{GitHub}
    \enspace$|$\enspace
  \href{https://seshat-ai.com}{Seshat-AI}
    \enspace$|$\enspace
  \href{https://kpoviessi-stephane.vercel.app/}{Portfolio}
\end{minipage}
\vspace{3pt}

%-----------PROFIL-----------
\section{Profil}
\small{
Ingénieur Big Data \& Data Science orienté AI Engineering, je conçois des systèmes agentiques combinant LLM, RAG, orchestration d’outils et automatisation. Développeur de Seshat, plateforme IA construite autour d’une architecture Python/Go, d’agents, de services backend et de composants orientés production. Expérience chez Allianz France en Data Engineering, automatisation des tests et fiabilisation de traitements à grande échelle.
}
\vspace{4pt}


%-----------EXPERIENCES-----------
\section{Expériences Professionnelles}
\resumeSubHeadingListStart

\resumeSubheading
{Ingénieur Data / Data Scientist}{Nov. 2025 -- Avr. 2026}
{Allianz France | Direction P4D, Équipe Advanced Data \& Climate}{Paris, France}
\resumeItemListStart
  \resumeItem{Migration d'un datamart Assurance de SAS vers \textbf{Python/PySpark}, avec validation de la parité fonctionnelle sur plus de \textbf{15 000 lignes de code}.}
  \resumeItem{Implémentation d'une architecture \textbf{Médaillon} sur Azure Databricks intégrant plus de \textbf{150 règles de gestion métier} pour fournir des datasets fiables aux cas d'usage analytiques.}
  \resumeItem{Conception et automatisation d'un pipeline de \textbf{Data Quality géospatiale} pour fiabiliser le reporting et les traitements aval.}
  \resumeItem{Développement d'un framework de \textbf{tests automatisés SAS/Python} pour la validation de non-régression et la préparation à la production.}
\resumeItemListEnd

\resumeSubheading
{Data Analyst -- Missions d'appui stratégique}{2023 -- 2025}
{LYNE ET FRERES SARL \normalfont\small(Missions récurrentes d'analyse Data)}{Ikpinlè, Bénin / À distance}
\resumeItemListStart
  \resumeItem{Analyse de données commerciales, opérationnelles et de marché avec \textbf{Python/Pandas} pour soutenir les décisions de diversification et le suivi des indicateurs clés de performance.}
  \resumeItem{Conception de rapports et tableaux de bord interactifs avec \textbf{Streamlit} et interfaces web HTML/CSS/JavaScript pour restituer les analyses au personnel de l'entreprise.}
  \resumeItem{Suivi analytique des activités agricoles, notamment les plantations de palmier et les étapes de transformation, afin d'identifier des pistes d'amélioration opérationnelle.}
  \resumeItem{Mise à disposition et accompagnement à l'usage responsable d'outils IA internes basés sur \textbf{OpenWebUI} et modèles open source pour appuyer les usages métiers.}
\resumeItemListEnd

\resumeSubheading
{Ingénieur Machine Learning}{Mai 2024 -- Août 2024}
{Groupe Sylvagreg \normalfont\small(Stage)}{Lille, France} 
\resumeItemListStart
  \resumeItem{Recherche et développement d'une solution de reconstruction 3D à partir d'images 2D, combinant Computer Vision, photogrammétrie et modèles génératifs sous PyTorch.}
  \resumeItem{Automatisation des pipelines de prétraitement, d'augmentation de données et de génération pour fiabiliser les traitements d'inférence.}
  \resumeItem{Déploiement cloud sur \textbf{AWS} avec conteneurisation \textbf{Docker} et interface de visualisation pour les utilisateurs métiers.}
\resumeItemListEnd

\resumeSubHeadingListEnd

%-----------PROJECTS-----------
\section{Projets Sélectionnés}
\resumeSubHeadingListStart

\resumeProjectHeading
{\textbf{Seshat \& Seshat-AI} $|$ \emph{Python, Go, TypeScript, Electron/React, RAG, MCP, PostgreSQL, pgvector, Chroma}}{\href{https://github.com/EngineerProjects/seshat}{GitHub}}
\resumeItemListStart
  \resumeItem{Conception et développement de bout en bout d'une \textbf{plateforme IA modulaire et industrialisable}, incluant services Python/Go, backend local embarqué, backend cloud multi-tenant, application desktop Electron/React et console d'administration.}
  \resumeItem{Construction d'un moteur \textbf{RAG hybride} : ingestion documentaire, chunking, embeddings multi-fournisseurs, indexation PostgreSQL/\textbf{pgvector} et \textbf{Chroma}, recherche vectorielle et restitution sourcée.}
  \resumeItem{Orchestration \textbf{multi-agents} capable d'exécuter des outils externes via le protocole \textbf{Model Context Protocol (MCP)}, avec automatisation de tâches et chaînage d'actions applicatives.}
  \resumeItem{Conception d'une couche d'abstraction \textbf{LLM multi-fournisseurs} (OpenAI, Mistral, Gemini, Ollama) permettant de partager des capacités communes entre génération de texte, embeddings, image et audio.}
  \resumeItem{Implémentation de fonctionnalités orientées entreprise : permissions, auditabilité, supervision \textbf{human-in-the-loop}, tâches planifiées et automatisation de workflows internes.}
  \resumeItem{Sécurisation applicative : stockage des identifiants via trousseaux natifs (DPAPI/Keychain/libsecret), authentification, gestion de sessions et isolation des espaces de travail.}
\resumeItemListEnd

\resumeProjectHeading
{\textbf{Tech Watch Agent} $|$ \emph{Python, FastAPI, LangGraph, PostgreSQL, pgvector, React, Docker}}{\href{https://github.com/EngineerProjects/tech-watch-agent}{GitHub}}
\resumeItemListStart
  \resumeItem{Développement d'une plateforme de \textbf{veille technologique multi-agents} automatisant la planification, la collecte multi-source, la synthèse d'informations et la génération de rapports ciblés.}
  \resumeItem{Backend FastAPI avec orchestration \textbf{LangGraph}, mémoire sémantique PostgreSQL/\textbf{pgvector}, configuration LLM/embeddings et industrialisation Docker avec frontend React.}
\resumeItemListEnd

\resumeProjectHeading
{\textbf{BI Retail Solution} $|$ \emph{Python, PySpark, Airflow, PostgreSQL, Power BI, Docker}}
{\href{https://github.com/EngineerProjects/BI_Retail}{GitHub}}
\resumeItemListStart
    \resumeItem{Mise en place d'une chaîne Data de bout en bout pour l'exploitation des données de ventes et de supply chain : génération et ingestion des données, transformations \textbf{PySpark} et chargement dans un \textbf{Data Warehouse PostgreSQL}.}
    \resumeItem{Orchestration des traitements avec \textbf{Apache Airflow} et Docker Compose, avec tableaux de bord \textbf{Power BI} couvrant les ventes, les clients, les stocks et la performance logistique.}
\resumeItemListEnd

\resumeProjectHeading
{\textbf{Robot autonome de collecte de déchets} $|$ \emph{Python, YOLOv8, OpenCV, Reinforcement Learning, Dijkstra}}
{AceObject}
\resumeItemListStart
  \resumeItem{Développement d'un pipeline de vision par ordinateur avec \textbf{YOLOv8} et \textbf{OpenCV} pour détecter, classifier et suivre les déchets tout en les distinguant des obstacles.}
  \resumeItem{Implémentation d'une planification de chemin basée sur l'algorithme de \textbf{Dijkstra} pour calculer des itinéraires de navigation efficaces pour la collecte autonome de déchets (environnements tels que les plages).}
  \resumeItem{Exploration de l'\textbf{apprentissage par renforcement} pour la prise de décision autonome et l'évitement d'obstacles, combinant perception visuelle et navigation robotique.}
\resumeItemListEnd

\resumeSubHeadingListEnd


%-----------SKILLS-----------
\section{Compétences}
\begin{itemize}[leftmargin=0.15in, label={}]
\small{\item{
\textbf{IA Générative \& Agents} : LLMs, RAG, Systèmes Multi-Agents, MCP, Tool Calling, LangGraph, LangChain, Human-in-the-Loop, Automatisation de Workflows IT \\
\textbf{Programmation \& Backend} : Python, Go, SQL, FastAPI, REST APIs, PostgreSQL, SSE \\
\textbf{Cloud \& Industrialisation} : Azure (ADLS, ADF, Databricks), AWS (EC2, S3), Docker, Git, CI/CD \\
\textbf{Data Engineering} : PySpark, Delta Lake, ETL/ELT, Apache Airflow, Architecture Médaillon, Data Quality \\
\textbf{Machine Learning} : Pandas, Scikit-learn, PyTorch, TensorFlow/Keras, Computer Vision, Forecasting, Évaluation de Modèles
}}
\end{itemize}


%-----------FORMATION-----------
\section{Formation}
\resumeSubHeadingListStart
  \resumeSubheading
    {JUNIA ISEN}{Lille, France}
    {Diplôme d'Ingénieur | Big Data \& Data Science}{Diplômé en 2026}
  \resumeSubheading
    {JUNIA ISEN}{Rabat, Maroc}
    {Licence | Électronique et Systèmes Numériques}{2023}
  \resumeSubheading
    {Cycle préparatoire Marie Stella}{Cotonou, Bénin}
    {Classe préparatoire MPSI | Mathématiques, Physique, Sciences de l'Ingénieur}{2022}
\resumeSubHeadingListEnd

%-----------CERTIFICATIONS-----------
\section{Certifications}
\begin{itemize}[leftmargin=0.15in, label={}]
\small{\item{
\textbf{Supervised Machine Learning} | \href{https://www.coursera.org/learn/machine-learning}{DeepLearning.AI} $|$ 2023 \vspace{1.5pt} \\
\textbf{Microsoft AI Engineer Career Path} | Microsoft Learn $|$ En cours
}}
\end{itemize}

%-----------LANGUES-----------
\section{Langues}
\begin{itemize}[leftmargin=0.15in, label={}]
  \small{\item{
   \textbf{Français} : Langue maternelle \vspace{1.5pt} \\
   \textbf{Anglais} : B2 | Linguaskill Business English \vspace{1.5pt} \\
   \textbf{Yoruba, Fon, Goun} : Langues maternelles (Bilingue)
  }}
\end{itemize}

%-----------ENGAGEMENTS-----------
\section{Engagements \& Associations}
\resumeSubHeadingListStart
  \resumeSubheading
    {Vice-Président | Commission Réseau JUNIA ALUMNI}{Sep. 2024 -- Août 2025}
    {Coordination du réseau, engagement des alumni et organisation d'événements professionnels}{Lille, France}

  \resumeSubheading
    {Responsable Communication | Commission Réseau JUNIA ALUMNI}{Sep. 2023 -- Août 2024}
    {Communication et promotion des initiatives alumni}{Lille, France}
\resumeSubHeadingListEnd

\end{document}
